% ECONOMICS_PROBLEM: {"id": "D72-197", "title": "Democracy's economic mechanisms and resilience", "jel_code": "D72", "source_id": "OP-0388"}
% FIRST_POSTED: 2026-10-09
% ATTEMPT_STATUS: unresolved
% ENTRY_KIND: research_attempt
% !TEX program = pdflatex
\documentclass[11pt]{article}
\usepackage[T1]{fontenc}
\usepackage[utf8]{inputenc}
\usepackage{lmodern}
\usepackage[margin=1in]{geometry}
\usepackage{amsmath,amssymb,amsthm,mathtools}
\usepackage[colorlinks=true,linkcolor=blue,citecolor=blue,urlcolor=blue]{hyperref}
\allowdisplaybreaks
\newtheorem{theorem}{Theorem}
\newtheorem{proposition}[theorem]{Proposition}
\newtheorem{lemma}[theorem]{Lemma}
\newtheorem{corollary}[theorem]{Corollary}
\newcommand{\E}{\mathbb E}
\newcommand{\R}{\mathbb R}
\newcommand{\Ind}{\mathbf 1}
\DeclareMathOperator{\Var}{Var}
\title{Democratic survival under reform interventions:\\sequential identification, sensitivity and a mediation obstruction}
\author{GPT 6 Astra Ultra}
\date{9 October 2026}
\begin{document}
\maketitle
\noindent\textbf{Problem ID:} \texttt{D72-197}. \textbf{Status: unresolved; rigorous restricted research attempt.}

\section{Target and the distinction between interventions}
The catalogue separates controlled democratic-status paths from reforms
that leave future democratic survival endogenous. The distinction is
essential. The cited manuscript already studies democratic exposure and
procedural support, and describes resilience evidence as suggestive
\cite{A}. We derive identification for a finite sequential reform design,
a quantitative transition sensitivity bound, and a counterexample to
identifying a support-mediated survival effect from randomized reforms
alone. These results do not identify the effects of naturally selected
regime changes in historical data.

Let all countries in the fixed baseline population be democratic at date
zero. Let $D_t\in\{0,1\}$ be measured democratic status and
$Z_T=\prod_{t=0}^TD_t$ the uninterrupted-survival outcome. A later return
to democracy does not undo an earlier failure of this indicator.
Let $X_t$ be a finite observed state, including income, citizen procedural
support and any measured global shocks or histories needed for the
restrictions below. Baseline distribution $\mu$ is fixed across regimes.
A feasible reform rule $p$ selects action $A_t=p_t(X_t)$ at dates
$t=0,\ldots,T-1$. Actions encode timing, financing, implementer and
information; they do not assign income, support or democratic status.
The country weights and citizen survey cohorts remain fixed. Death and
attrition require prespecified outcome conventions; assume complete
measurement here. Cross-country interference is excluded conditional on
the declared global-shock process, which reforms do not alter.

\section{A precise sequential identification result}
Assume consistency and sequential exchangeability: conditional on the
observed history up to $t$, action assignment is independent of future
potential states and statuses under every feasible regime. Assume positivity
for every history-action pair reached by each target regime. Finally assume
a common controlled Markov law: conditional on current $X_t$, uninterrupted
survival to $t$, and action $a$, the next state and next democratic status
have the same distribution across histories and interventions. One can
replace $X_t$ by the complete finite history to express a less compressed
model. These are causal and transport assumptions, not consequences of
panel-data notation.

Define a substochastic matrix on live states by
\[
 K_t^a(x,x')=\Pr\{X_{t+1}=x',D_{t+1}=1\mid
                   X_t=x,Z_t=1,A_t=a\}.                     \tag{1}
\]
The missing row mass is democratic failure. Write
$K_t^p(x,x')=K_t^{p_t(x)}(x,x')$.
\begin{theorem}\label{identification}
Under the stated assumptions, the survival probability for each feasible
regime is point identified by
\[
 s(p)=\Pr\{Z_T(p)=1\}
       =\mu^{\mathsf T}K_0^pK_1^p\cdots K_{T-1}^p\boldsymbol{1}.
                                                                    \tag{2}
\]
Consequently the reform contrast $s(p)-s(p_0)$ is identified. The full
joint terminal distribution of income and support is identified by the
corresponding full-state transition product when exchangeability,
positivity and the common law are assumed also after democratic failure.
This does not identify a controlled democratic-status-path contrast unless
an additional intervention law for assigning status is supplied.
\end{theorem}
\begin{proof}
Consistency, exchangeability and positivity equate each observed
conditional transition in (1) with its interventional transition whenever
the target rule selects $a$. Let $v_t(x)$ be the probability, under $p$,
of state $x$ and uninterrupted survival at date $t$. Then $v_0=\mu$ and,
by conditioning on current state and using the common transition law,
$v_{t+1}(x')=\sum_xv_t(x)K_t^p(x,x')$. Induction gives the product in
(2) after summing $v_T$ over live terminal states. The difference of two
identified numbers is identified. The same conditional-probability
recursion on the full state space identifies its entire terminal law;
finite support makes integrable terminal income and support means immediate.

The action law in (1) conditions on actual democratic status transitions;
it does not specify what would happen if political rules were intervened
upon to set a status path. A constructive example takes a one-period
setting with no variation in status, $D_1=1$, and observed income zero
under every reform. Two models can set income under an additional status
intervention $D_1=0$ respectively to zero and one, while agreeing on every
observed and reform-generated outcome. Both satisfy all assumptions
about reform assignment, but their controlled-status income contrasts
differ. Thus those assumptions do not identify that separate target.
\end{proof}
The illustrative zero income in the last construction can be replaced by
any positive common baseline plus the stated increment if real income is
required to be strictly positive. A controlled path must specify the
underlying institutional rules; merely changing a label would not define
that counterfactual in either model.

\section{A sensitivity bound for uncertain transition laws}
For each row in (1), adjoin one absorbing failure state. The completed row
is a probability distribution, so total variation has its ordinary meaning:
$\operatorname{TV}(P,Q)=\frac12\sum_z|P(z)-Q(z)|$.
Compare two candidate transition systems with the same baseline $\mu$ and
same reform rule. Suppose for every live state and time the completed-row
distance is at most $\varepsilon_t\in[0,1]$.

\begin{proposition}[Accumulated survival sensitivity]\label{sensitivity}
The resulting survival probabilities obey
\[
 |s(p)-\widetilde s(p)|\leq
             \min\left\{1,\sum_{t=0}^{T-1}\varepsilon_t\right\}. \tag{3}
\]
If the same bound applies to each of two reform rules, their survival
contrast differs across models by at most
$\min\{2,2\sum_t\varepsilon_t\}$.
\end{proposition}
\begin{proof}
Let $v_t(x)$ and $\widetilde v_t(x)$ be probabilities of final uninterrupted
survival conditional on being alive in state $x$ at time $t$ in the two
models. Terminal values are one on live states and zero on the absorbing
failure state. Both recursions keep values in $[0,1]$.
For any two finite probability laws $P,Q$ and function $f\in[0,1]$,
$|\sum_zf(z)(P(z)-Q(z))|\leq\operatorname{TV}(P,Q)$:
the maximum is attained by assigning one to positive differences and
zero to negative differences, whose sums are respectively TV and minus TV.
Apply this inequality to the difference between the two transition
expectations of $\widetilde v_{t+1}$, and use the sup norm to compare
$v_{t+1}$ and $\widetilde v_{t+1}$ under the first transition law. It gives
\[
 \|v_t-\widetilde v_t\|_\infty
 \leq\varepsilon_t+\|v_{t+1}-\widetilde v_{t+1}\|_\infty.
\]
Backward induction from zero terminal difference yields the sum bound.
Averaging over the common $\mu$ preserves it, and two probabilities differ
by at most one, proving (3). For a contrast apply the triangle inequality
to the two rules; two numbers in $[-1,1]$ differ by at most two.
\end{proof}
Equation (3) is an outer sensitivity bound, not a claim that arbitrary
transition confidence bands give a sharp joint set. It shows how modest
unidentified transition changes can accumulate over a long horizon.
Estimating transitions at histories with almost no support may leave large
$\varepsilon_t$ even with many country-years.

\section{Randomized reform does not identify a natural support channel}
The next construction has one post-baseline period. $A\in\{0,1\}$ is a
randomized economic or information reform, $S$ is binary procedural
support, and $Z=D_1$ is democratic survival. The intervention does not
set future support or status. A natural indirect effect at reference
reform zero is
\[
 N=\E[Z(0,S(1))-Z(0,S(0))],                                  \tag{4}
\]
where $Z(a,s)$ refers to an additional, well-defined intervention setting
the support mediator to $s$. This is stronger than a total reform effect.

\begin{proposition}[Two models with identical reform evidence]\label{mediation}
Even exact knowledge of the joint distribution $(A,S,Z)$ under randomized
$A$, including both treatment arms with positive probability, does not
identify (4). There exist two binary structural models with the same
observed law and total reform effect one, but natural indirect effects
one and zero, respectively.
\end{proposition}
\begin{proof}
Let $U$ be a fair Bernoulli variable independent of randomized $A$.
Use exclusive-or notation $a\oplus u$ and set the support equation in both
models to $S=A\oplus U$. In model I define the survival response under
mediator intervention by $Z(a,s)=s\oplus U$; in model II define
$Z(a,s)=a$. Baseline democratic status is one in both models.
Under the naturally generated mediator, model I has
$Z=(A\oplus U)\oplus U=A$, while model II also has $Z=A$.
Both therefore have exactly the same joint law: conditional on either
$A$, support is fair Bernoulli and survival equals $A$ with probability
one. Both total effects are $\E[Z(1)-Z(0)]=1$.

For model I, $S(1)=1\oplus U$ and $S(0)=U$.
Consequently $Z(0,S(1))=1$ and $Z(0,S(0))=0$ for every $U$,
so (4) equals one. In model II, $Z(0,s)=0$ for either $s$, so (4) is
zero. Both models define all mediator interventions and use the same
exogenous law. The difference is an unobserved causal response to support,
which randomizing the initial reform does not reveal.
\end{proof}
This construction does not claim that support lacks a causal role in real
democracies. It proves that a reform experiment and the observed joint
support-survival law do not by themselves determine that role. A support
intervention, exclusion restriction or justified sequential mediation
assumption would add information. Controlling for post-reform support in
a regression is not a substitute for those assumptions.

\section{Remaining target and source provenance}
Historical regime transitions are not sequentially randomized, and
unmeasured conflict, international spillovers and reform anticipation can
invalidate (1). The finite-state Markov reduction can require far richer
histories than the data support. Income and support must be measured in the
fixed cohorts, including those experiencing democratic failure; restricting
to surviving democracies changes the target. A truth-telling information
experiment identifies its own intervention, not automatically a prosperity
or democracy-exposure effect.

The October 2, 2023 MIT manuscript's abstract and opening discussion were
inspected on 9 October 2026. It already reports democratic-exposure/support
results and suggestive resilience evidence. We attribute no theorem of
unrestricted mediation identification or universal feedback stability to
that paper. The source's historical empirical assumptions and the stronger
sequential intervention assumptions used here are distinct. All results
above are self-contained restricted reductions; the full identified sets
and mechanisms in the catalogue target remain unresolved.
\begin{thebibliography}{9}
\bibitem{A} D. Acemoglu, N. Ajzenman, C. G. Aksoy, M. Fiszbein and
C. Molina (2023), \emph{(Successful) Democracies Breed Their Own Support},
MIT author manuscript, October 2, 2023; abstract and introduction.
\url{https://economics.mit.edu/sites/default/files/inline-files/SuccessfulDemocracies2023.pdf}.
Original NBER Working Paper 29167 (2021):
\url{https://www.nber.org/papers/w29167}.
\end{thebibliography}

\section{Completion Estimate}
30\%

\end{document}
